\documentclass[11pt]{article}
\usepackage{ochamath}
\subjectcolor{2F5DA8}
\setsheet{線形代数}{二次形式と正定値　演習}{https://math.ochanote.com/linear-algebra/quadratic-forms/}
\namelinetrue
\begin{document}
\makesheethead{問題}

\problem[展開と平方完成]
$A=\begin{pmatrix}3&1\\1&2\end{pmatrix}$ の二次形式を展開し、正定値であることを平方完成で示せ。
\workspace{38mm}

\problem[主軸変換]
$A=\begin{pmatrix}2&1\\1&2\end{pmatrix}$ を使う二次形式を、正規直交な固有基底の座標 $z_1,z_2$ で表せ。
\workspace{38mm}
\newpage

\problem[符号の分類]
$\operatorname{diag}(1,0)$ と $\operatorname{diag}(1,-1)$ を、正定値・半正定値・不定値の観点で分類せよ。
\workspace{38mm}

\problem[ばねのエネルギー]
$K=\begin{pmatrix}2&-1\\-1&2\end{pmatrix}\ \mathrm{N/m}$、変位 $\boldsymbol x=(0.01,0)^{\mathsf T}\ \mathrm m$ のエネルギー $U=\frac12\boldsymbol x^{\mathsf T}K\boldsymbol x$ を求めよ。
\workspace{38mm}
\end{document}
